\scanpage{300}
\mspage{- 160 -}\hfill IX

d'entiers $[1,r]$), $D_\lambda$ étant l'image par $\underline{f}:\underline{X}\handwritten{\longrightarrow}\underline{S}$ de la composante $D'_\lambda$ passant par le point singulier $x_\lambda$ de $X_{\unsure{o|0}}$ du lieu de non lissité $D'=\handwritten{\bigcup}_{\lambda\in L}D'_\lambda$ du morphisme $\underline{f}$. Par construction même, $f:X\handwritten{\longrightarrow} S$ est isomorphe à l'image inverse de $\underline{f}:\underline{X}\handwritten{\longrightarrow}\underline{S}$ par un $W$-morphisme convenable $g:S\handwritten{\longrightarrow}\underline{S}$ :

\diagram{\[
\begin{tikzcd}[ampersand replacement=\&]
 \underline{X} \arrow[d,"\underline{f}"'] \& X \arrow[l] \arrow[d,"f"]\\
 \underline{S} \& S \arrow[l,"g"'] \& .
\end{tikzcd}\tag*{(12.8.2)}
\]}

Posons
\[
 P=\underline{\mathrm{Pic}}^{o}_{\underline{X}/\underline{S}}\qquad ;\tag*{(12.8.3)}
\]
comme la fibre spéciale de $\underline{f}$ est géométriquement réduite de dimension 1, il s'ensuit que $\underline{F}$ est cohomologiquement plat en dimension 0 (EGA III 7.8.1) et que $P$ \deleted{\unsure{$-$}}est formellement lisse sur $S$ (ce dernier fait résultant de ce que $H^2(X_o,\underline{O}_{X_o})=0$, $X_o$ étant de dimension 1). Pour simplifier, nous allons admettre que $P$ \uline{est représentable par un schéma en groupes sur} $S$ (qui sera nécessairement un schéma en groupes lisse à fibres connexes) ; c'est là un cas particulier d'un résultat inédit de MUMFORD cité dans\handwritten{\unsure{\textsuperscript{1}}} [22](considérablement généralisé par RAYNAUD dans la note citée, dans le cas où la base est un trait). Comme aucune démonstration des résultats de MUMFORD ni de RAYNAUD n'a été publiée, le lecteur préfèrera peut-être utiliser les résultats généraux de M. ARTIN [1], qui donnent la représentabilité de $P$ (sous la seule hypothèse de partir d'un $\underline{f}$ propre, plat et cohomologiquement plat en dimension 0) par un ``espace algébrique''\unsure{.}

\mspage{478}
% archive sheet number: 301
